%%% FIELD kink-residual-statement
There is a finite constant \(K_{\mathrm{res}}\), depending only on \((c,C,G)\), such that, for every \(N\geq N_0\), on the scaled lattice \(E_{1,N}\),
\[
 \sup_{p\in\mathcal P,\,|\gamma|\leq G}
 \mathbb E_{\rho_{1,N}}
 \left|\{G_{1,N}-\mathcal L_{\gamma_N}\}
 h_{1,N}(X_{1,N})\right|
 \leq K_{\mathrm{res}}N^{-1/2}.
\]
The inequality includes the physical boundary states and every state whose available upward or downward edge meets or crosses \(\gamma_N\).
%%% FIELD kink-residual-proof
Fix \(N\geq N_0\), \(p\in\mathcal P\), and \(|\gamma|\leq G\).  Put
\[
 \delta=N^{-1/2},\qquad g=\gamma_N,
 \qquad b_g(x)=-dx-\nu\min\{x,g\},
 \qquad \overline G=G+1.
 \tag{1}
\]
By the definition of \(M_N\),
\[
 0\leq g-\gamma<\delta\leq1,
 \qquad |g|\leq\overline G.
 \tag{2}
\]
The compact rate assumptions give the strictly positive, uniform bounds
\[
 2c\leq d\leq2C,\qquad 3c\leq D\leq3C,\qquad
 q_*:=\frac{2c^2}{3C}\leq q=\frac{\lambda(\tau+\nu)}D\leq C.
 \tag{3}
\]
Also,
\[
 a(p)\geq\frac{c}{3C}\geq\beta,
 \qquad 1-a(p)\geq\frac{2c}{C+2c}\geq\beta.
\]
Writing \(r_0=(G+\sqrt{G^2+4\beta})/(2\beta)\), the definition of \(N_0\) gives \(\sqrt N\geq r_0\), and \(r_0\) is the positive root of \(\beta r^2-Gr-1=0\).  The quadratic is increasing on \([r_0,\infty)\), so
\[
 \beta N-G\sqrt N\geq1,
 \qquad
 1\leq M_N=\lceil Na(p)+\gamma\sqrt N\rceil\leq N-1.
\]
Thus the displayed birth and death rates are nonnegative, and their only unavailable boundary jumps have rate zero.  Every division by \(d,D,q\) below is valid.  The relevant state support is the finite lattice \(E_{1,N}=\{(k-Na)/\sqrt N:0\leq k\leq N\}\).  No rank, overlap, or limiting argument is used in this finite-\(N\) estimate.

We first obtain bounds for the diffusion Poisson solution uniformly for \(|g|\leq\overline G\).  The density in the piecewise Poisson construction satisfies
\[
 q\pi_g'(x)=b_g(x)\pi_g(x)
 \tag{4}
\]
on both sides of \(g\), and the density and its first derivative agree at \(g\), because \(b_g\) is continuous there.  Its displayed Gaussian branches have finite moments of every fixed order.  Hence integration by parts, with the Gaussian boundary terms equal to zero, gives
\[
 0=\int \mathcal L_g(x^2)\pi_g(x)\,dx
   =-2d\mathbb E_{\pi_g}X^2
     -2\nu\mathbb E_{\pi_g}\{X\min(X,g)\}+2q.
 \tag{5}
\]
For every real \(x,g\),
\[
 x\min(x,g)\geq-|g|\,|x|.
 \tag{6}
\]
Indeed, the left side is nonnegative when \(x\leq0\) or \(g\geq0\), except possibly when \(x>0>g\), in which case it equals \(xg=-|g||x|\).  If \(y_g=(\mathbb E_{\pi_g}X^2)^{1/2}\), equations (2), (3), (5), and (6), followed by Cauchy--Schwarz, imply
\[
 2c y_g^2\leq C+C\overline G y_g.
 \tag{7}
\]
Consequently
\[
 |m_g|\leq y_g\leq
 M_\pi:=\frac{C\overline G+\sqrt{C^2\overline G^2+8cC}}{4c}.
 \tag{8}
\]

For completeness, the global first-derivative bound used below follows directly from the primitive drift.  Couple two solutions of
\[
 dY_t=b_g(Y_t)\,dt+\sqrt{2q}\,dB_t
 \tag{9}
\]
with the same Brownian motion and initial states \(y>x\).  Since \(u\mapsto\min(u,g)\) is nondecreasing and one-Lipschitz,
\[
 -D(Y_t^y-Y_t^x)
 \leq b_g(Y_t^y)-b_g(Y_t^x)
 \leq-d(Y_t^y-Y_t^x).
 \tag{10}
\]
Until a first meeting time, the Brownian terms in (9) cancel, so the positive difference is absolutely continuous and satisfies (10) almost everywhere.  Gronwall's inequality gives the following bounds up to that time; the strictly positive lower bound rules out a finite meeting, so they hold for every \(t\geq0\):
\[
 e^{-Dt}(y-x)\leq Y_t^y-Y_t^x\leq e^{-dt}(y-x).
 \tag{11}
\]
Coupling one initial state to an independent stationary draw from \(\pi_g\) shows that
\[
 r_g(x):=\int_0^\infty\{\mathbb E_xY_t-m_g\}\,dt
 \tag{12}
\]
is absolutely convergent.  The semigroup identity
\[
 P_s r_g-r_g=-\int_0^sP_t(x-m_g)\,dt
 \tag{13}
\]
and dominated convergence yield \(-\mathcal L_g r_g=x-m_g\).  Fubini's theorem and stationarity give \(\int r_g\,d\pi_g=0\); absolute integrability follows from (11) and the finite first moment established above.  The centered solution is unique: the derivative of the difference of two solutions is a constant times \(1/\pi_g\) by (4), and the Gaussian tails make such a difference nonintegrable unless that constant is zero; centering then removes the remaining additive constant.  Thus (12) is the function \(h_g\) in the diffusion Poisson construction.  Integrating (11) over \(t\) and taking difference quotients therefore gives
\[
 \frac1D\leq h_g'(x)\leq\frac1d\leq\frac1{2c},
 \qquad x\in\mathbb R.
 \tag{14}
\]
This also independently verifies, on the one-step-enlarged window required by the ceiling in (2), the derivative control appearing in Theorem \(\mathrm{diffusion\mbox{-}poisson}\).

Let
\[
 F_g(x)=\int_{-\infty}^x(u-m_g)\pi_g(u)\,du.
 \tag{15}
\]
The construction \(h_g'=-F_g/(q\pi_g)\), equation (4), and positivity of \(q\pi_g\) show that \(h_g'\) is continuously differentiable.  The Poisson equation can therefore be written pointwise as
\[
 qh_g''(x)=-(x-m_g)-b_g(x)h_g'(x).
 \tag{16}
\]
Because
\[
 |b_g(x)|\leq3C|x|+C\overline G,
 \tag{17}
\]
equations (3), (8), (14), and (16) give
\[
 |h_g''(x)|\leq C_1(1+|x|),
 \quad
 C_1:=q_*^{-1}\left\{1+M_\pi+\frac{3C+C\overline G}{2c}\right\}.
 \tag{18}
\]
On each open side of \(g\), differentiating (16) gives
\[
 qh_g'''(x)=-1-b_g'(x)h_g'(x)-b_g(x)h_g''(x),
 \qquad b_g'(x)\in\{-D,-d\}.
 \tag{19}
\]
It follows from (3), (14), (17), and (18) that, for almost every \(x\),
\[
 |h_g'''(x)|\leq C_2(1+x^2),
 \quad
 C_2:=q_*^{-1}\left\{1+\frac{3C}{2c}
       +2C_1(3C+C\overline G)\right\}.
 \tag{20}
\]
Both \(b_g\) and \(h_g'\) are continuous at \(g\), so (16) makes \(h_g''\) continuous there.  Since (19)--(20) hold on the two sides, \(h_g''\) is locally absolutely continuous across \(g\).  Taylor's formula with integral remainder is therefore valid on an interval that meets or crosses the kink, with the essential supremum in (20).  This is the step that prevents a kink-edge remainder of order one.

At \(x=(k-Na)/\sqrt N\), define the available birth and death rates
\[
 B_N(x)=\lambda(N-k),
 \qquad D_N^-(x)=\tau k+\nu\min\{k,M_N\}.
 \tag{21}
\]
Since \(k=Na+\sqrt N x\), \(M_N=Na+\sqrt N g\), and \(\lambda(1-a)=(\tau+\nu)a=q\), direct algebra gives the exact identities
\[
 \delta\{B_N(x)-D_N^-(x)\}=b_g(x),
 \tag{22}
\]
\[
 \frac{\delta^2}{2}\{B_N(x)+D_N^-(x)\}
 =q+\frac\delta2 r_g^N(x),
 \qquad
 r_g^N(x)=(\tau-\lambda)x+\nu\min(x,g).
 \tag{23}
\]
These identities also hold at \(k=0,N\), because the unavailable rate in (21) is then exactly zero.  Moreover,
\[
 B_N(x)+D_N^-(x)\leq3CN,
 \qquad
 |r_g^N(x)|\leq C_r(1+|x|),
 \quad C_r:=C(2+\overline G).
 \tag{24}
\]

Apply Taylor's formula only to available edges.  Equations (18) and (20) yield
\[
 h_g(x\mathbin{\pm}\delta)
 =h_g(x)\mathbin{\pm}\delta h_g'(x)
   +\frac{\delta^2}{2}h_g''(x)+\varepsilon_\pm(x),
 \tag{25}
\]
where
\[
 |\varepsilon_\pm(x)|
 \leq\frac{\delta^3C_2}{6}
       \sup_{|u-x|\leq\delta}(1+u^2)
 \leq\frac{\delta^3C_2}{2}(1+x^2).
 \tag{26}
\]
Substituting (22)--(26) into the scaled lattice generator gives
\[
 G_{1,N}h_g(x)
 =b_g(x)h_g'(x)+qh_g''(x)
  +\frac\delta2r_g^N(x)h_g''(x)+R_N(x),
 \tag{27}
\]
with
\[
 |R_N(x)|
 \leq\{B_N(x)+D_N^-(x)\}
       \max\{|\varepsilon_+(x)|,|\varepsilon_-(x)|\}
 \leq\frac{3CC_2}{2}\delta(1+x^2).
 \tag{28}
\]
Since \((1+|x|)^2\leq2(1+x^2)\), equations (18), (24), (27), and the definition of \(\mathcal L_g\) imply the pointwise, compact-uniform bound
\[
 |(G_{1,N}-\mathcal L_g)h_g(x)|
 \leq K_1\delta(1+x^2),
 \qquad
 K_1:=C_rC_1+\frac{3CC_2}{2}.
 \tag{29}
\]
Notice also that \(g=(M_N-Na)/\sqrt N\) is itself the lattice point corresponding to \(k=M_N\).  Thus a lattice edge can only remain on one smooth branch or have the kink as an endpoint.  Even without using this alignment, the absolute-continuity argument preceding (21) proves (25)--(29) for a genuinely crossing interval.

It remains to average (29) under the finite stationary law.  Applying the finite-state generator to \(x^2\) and using (22)--(23) gives exactly
\[
 G_{1,N}x^2=2xb_g(x)+2q+\delta r_g^N(x).
 \tag{30}
\]
By (2), (3), (6), and the sharper direct bound
\[
 |r_g^N(x)|\leq2C|x|+C\overline G,
 \tag{31}
\]
we have
\[
 G_{1,N}x^2
 \leq-2dx^2+L|x|+2C+C\overline G,
 \qquad L:=2C(\overline G+1).
 \tag{32}
\]
Young's inequality \(L|x|\leq dx^2+L^2/(4d)\), together with \(d\geq2c\), yields
\[
 G_{1,N}x^2\leq-dx^2+C_6,
 \qquad C_6:=2C+C\overline G+\frac{L^2}{8c}.
 \tag{33}
\]
The law \(\rho_{1,N}\) is stationary by definition and the treatment chain has the stated generator, so finite support permits termwise summation and gives \(\mathbb E_{\rho_{1,N}}G_{1,N}x^2=0\).  Therefore
\[
 \mathbb E_{\rho_{1,N}}X_{1,N}^2
 \leq\frac{C_6}{d}\leq\frac{C_6}{2c}.
 \tag{34}
\]
Taking stationary expectations in (29), using the lattice restriction \(h_{1,N}=h_g|_{E_{1,N}}\), and applying (34) proves
\[
 \mathbb E_{\rho_{1,N}}
 \left|\{G_{1,N}-\mathcal L_{\gamma_N}\}
 h_{1,N}(X_{1,N})\right|
 \leq K_1\left(1+\frac{C_6}{2c}\right)N^{-1/2}.
 \tag{35}
\]
Thus the required choice
\[
 K_{\mathrm{res}}=K_1\left(1+\frac{C_6}{2c}\right)
 \tag{36}
\]
is finite and depends only on \((c,C,G)\).  Since every bound before taking the supremum is uniform over \(p\in\mathcal P\) and \(|\gamma|\leq G\), equation (35) proves the asserted supremum inequality.

The estimate has the required centering consequence.  Stationarity gives \(\mathbb E_{\rho_{1,N}}G_{1,N}h_{1,N}=0\), while the Poisson equation gives \(-\mathcal L_g h_g=x-m_g\).  Hence
\[
 \left|\mathbb E_{\rho_{1,N}}X_{1,N}-m_{\gamma_N}\right|
 =\left|\mathbb E_{\rho_{1,N}}
   [(G_{1,N}-\mathcal L_{\gamma_N})h_{1,N}(X_{1,N})]\right|
 \leq K_{\mathrm{res}}N^{-1/2}.
 \tag{37}
\]
Equivalently,
\[
 \left|\mathbb E_{\rho_{1,N}}\frac{K_{1,N}}N
 -\left\{a(p)+\frac{m_{\gamma_N}}{\sqrt N}\right\}\right|
 \leq\frac{K_{\mathrm{res}}}{N}.
 \tag{38}
\]
This proves the order required for diffusion fitted centering without assuming it.

As an ancillary early-kill stress test, not used in the proof, take \(c=1/2\), \(C=2\), and \(G=2\), for which the stated formula gives \(N_0=600\).  We evaluated the stationary weights from the exact product
\[
 \rho_{1,N}(k)\ \propto\
 \prod_{j=0}^{k-1}
 \frac{\lambda(N-j)}{\tau(j+1)+\nu\min\{j+1,M_N\}}
 \tag{39}
\]
at all eight corners of the rate box, at thirty-nine nominal offsets between \(-1.9\) and \(1.9\), and at ceiling offsets \(0,0.11,0.37,0.63,0.89\).  Analytic piecewise-Gaussian formulas for \(h_g'\), integrated on each lattice edge by twenty-point Gauss--Legendre quadrature, gave maxima of \(\sqrt N\,\mathbb E_{\rho_{1,N}}|(G_{1,N}-\mathcal L_g)h_g|\) equal to \(0.0484345\), \(0.0479274\), and \(0.0475845\) at \(N=600,1200,2400\), respectively.  The worst corner was \((\lambda,\tau,\nu)=(1/2,1/2,2)\) near \(g=-0.2\).  This finite computation is consistent with (35), while (1)--(36), rather than the computation, certify every corner and ceiling phase.
